\documentclass{article}

% if you need to pass options to natbib, use, e.g.:
%     \PassOptionsToPackage{numbers, compress}{natbib}
% before loading neurips_2026

% The authors should use one of these tracks.
% Before accepting by the NeurIPS conference, select one of the options below.
% 0. "default" for submission
\usepackage{neurips_2026}
% the "default" option is equal to the "main" option, which is used for the Main Track with double-blind reviewing.
% 1. "main" option is used for the Main Track
%  \usepackage[main]{neurips_2026}
% 2. "position" option is used for the Position Paper Track
%  \usepackage[position]{neurips_2026}
% 3. "eandd" option is used for the Evaluations & Datasets Track
 % \usepackage[eandd]{neurips_2026}
% 4. "creativeai" option is used for the Creative AI Track
%  \usepackage[creativeai]{neurips_2026}
% 5. "sglblindworkshop" option is used for the Workshop with single-blind reviewing
 % \usepackage[sglblindworkshop]{neurips_2026}
% 6. "dblblindworkshop" option is used for the Workshop with double-blind reviewing
%  \usepackage[dblblindworkshop]{neurips_2026}
% 7. "education" option is used for the Education Track
%  \usepackage[education]{neurips_2026}


% After being accepted, the authors should add "final" behind the track to compile a camera-ready version.
% 1. Main Track
 % \usepackage[main, final]{neurips_2026}
% 2. Position Paper Track
%  \usepackage[position, final]{neurips_2026}
% 3. Evaluations & Datasets Track
 % \usepackage[eandd, final]{neurips_2026}
% 4. Creative AI Track
%  \usepackage[creativeai, final]{neurips_2026}
% 5. Workshop with single-blind reviewing
%  \usepackage[sglblindworkshop, final]{neurips_2026}
% 6. Workshop with double-blind reviewing
%  \usepackage[dblblindworkshop, final]{neurips_2026}
% Note. For the workshop paper template, both \title{} and \workshoptitle{} are required, with the former indicating the paper title shown in the title and the latter indicating the workshop title displayed in the footnote.
% For workshops (5., 6.), the authors should add the name of the workshop, "\workshoptitle" command is used to set the workshop title.
% \workshoptitle{WORKSHOP TITLE}
% 7. Education Track
% \usepackage[education, final]{neurips_2026}

% "preprint" option is used for arXiv or other preprint submissions
 % \usepackage[preprint]{neurips_2026}

% to avoid loading the natbib package, add option nonatbib:
%    \usepackage[nonatbib]{neurips_2026}

\usepackage[utf8]{inputenc} % allow utf-8 input
\usepackage[T1]{fontenc}    % use 8-bit T1 fonts
\usepackage{hyperref}       % hyperlinks
\usepackage{url}            % simple URL typesetting
\usepackage{booktabs}       % professional-quality tables
\usepackage{amsfonts}       % blackboard math symbols
\usepackage{nicefrac}       % compact symbols for 1/2, etc.
\usepackage{microtype}      % microtypography
\usepackage{xcolor}         % colors
\usepackage{pifont}
\usepackage{multirow}
\usepackage{amsmath,bm}

\usepackage{CJKutf8}
\newcommand{\red}[1]{{\color{red}#1}}
\newcommand{\todo}[1]{{\color{red}#1}}
\newcommand{\TODO}[1]{\textbf{\color{red}[TODO: #1]}}
% \newcommand{\ygm}[1]{\textcolor{purpler}{#1}}
\DeclareRobustCommand{\ygm}[1]{\textnormal{\textcolor{purple}{\begin{CJK}{UTF8}
{gbsn}#1\end{CJK}}}}
\DeclareRobustCommand{\words}[1]{\textnormal{\textcolor{blue}{\begin{CJK}{UTF8}
{gbsn}#1\end{CJK}}}}

% Note. For the workshop paper template, both \title{} and \workshoptitle{} are required, with the former indicating the paper title shown in the title and the latter indicating the workshop title displayed in the footnote. 
\title{Visual Instruction Tuning}


% The \author macro works with any number of authors. There are two commands
% used to separate the names and addresses of multiple authors: \And and \AND.
%
% Using \And between authors leaves it to LaTeX to determine where to break the
% lines. Using \AND forces a line break at that point. So, if LaTeX puts 3 of 4
% authors names on the first line, and the last on the second line, try using
% \AND instead of \And before the third author name.


\author{%
  David S.~Hippocampus\thanks{Use footnote for providing further information
    about author (webpage, alternative address)---\emph{not} for acknowledging
    funding agencies.} \\
  Department of Computer Science\\
  Cranberry-Lemon University\\
  Pittsburgh, PA 15213 \\
  \texttt{hippo@cs.cranberry-lemon.edu} \\
  % examples of more authors
  % \And
  % Coauthor \\
  % Affiliation \\
  % Address \\
  % \texttt{email} \\
  % \AND
  % Coauthor \\
  % Affiliation \\
  % Address \\
  % \texttt{email} \\
  % \And
  % Coauthor \\
  % Affiliation \\
  % Address \\
  % \texttt{email} \\
  % \And
  % Coauthor \\
  % Affiliation \\
  % Address \\
  % \texttt{email} \\
}


\begin{document}


\maketitle


\begin{abstract}
  Instruction tuning large language models (LLMs) using machine-generated instruction-following data has been shown to improve zero-shot capabilities on new tasks, but the idea is \textbf{less explored} in the \underline{multimodal field}.\ygm{(开篇点明任务背景:指令微调可以提升LLM零样本学习能力,但在多模态领域尚未得到充分探索)}
  We present the \textbf{first attempt} to use language-only GPT-4 to generate multimodal language-image instruction-following data.\ygm{(承接less explored,延伸创新点)}
  By instruction tuning on such generated data, we introduce LLaVA: \textbf{L}arge \textbf{L}anguage \textbf{a}nd \textbf{V}ision \textbf{A}ssistant, an end-to-end trained large multimodal model that connects a vision encoder and an LLM for general-prupose visual and language understanding.\ygm{(\ding{202}提出核心方法)}
  To facilitate future research on visual instruction following, we construct two evaluation benchmarks with diverse and challenging application-oriented tasks.\ygm{(\ding{203}构建评测bench)}
  Our experiments show that LLaVA demonstrates impressive multimodal chat abilities, sometimes exhibiting the behaviors of multimodal GPT-4 on unseen images/instructions, and yeilds a 85.1\% relative score compared with GPT-4 on a synthetic multimodal instruction-following dataset.\ygm{(\ding{204}总结实验效果)}
  When fine-tuned on Science QA, the synergy of LLaVA and GPT-4 achieves a new state-of-the-art accuracy of 92.53\%.\ygm{(\ding{205}易用性与泛化性证明)}
  We make GPT-4 generated visual instruction data, our model, and code publicly avaliable.
\end{abstract}


\section{Introduction}

Humans interact with the world through many channels such as vision and language, \textbf{as} each individual channel has a unique advantage in representing and communicating certain concepts, and thus facilitates a better understanding of the world.\ygm{(通过人类的多通道开篇)}
One of the core aspirations in artificial intelligence is to develop a general-purpose assistant that can effectively follow multi-model vision-and-language instructions, aligned with human intent to complete various real-world tasks in the wild.

To this end, the community has witnessed an emergent interest in developing language-augmented foundation vision models, with strong capabilities on open-world visual understanding such as classification, detection, segmentation, and captioning, as well as visual generation and editing. We refer readers to the \emph{Computer Vision in the Wild} reading list for a more up-to-date literature compilation.
In this line of work, each task is solved independently by one single large vision model, with the task instruction implicitly considered in the model design.
Future, language is only utilized to describe the image content.
While this allows language to play an important role in mapping visual signals to language semantics --- a common channel for human communication, it leads to models that usually have a fixed interface with interactivity and adaptability to the user's instruction.

Large Language Models (LLMs), on the other hand, have shown that language can play a wider role: a universal interface for a language-purpose assistant, where various task instructions can be explicitly represented in language and guide the end-to-end trained neural assistant to switch to the task of interest to solve it. For example, the recent success of ChatGPT and GPT-4 have demonstrated the power of aligned LLMs in following  human instructions, and have stimulated tremendous interest in developing open-source LLMs. Among them, LLaMA is an open source LLM that matches the performance of GPT-3. Alpaca, Vicuna, GPT-4-LLM utilize various machine-generated high quality instruction-following samples to improve the LLM's alignment ability, reporting impressive performance compared with proprietary LLMs. Importantly, this line of work is \emph{text-only}.

In this paper, we present visual instruction tuning, the first attempt to extend instruction-tuning to the language-image multimodal space, to pave the way towards building a general-purpose visual assistant. In patticular, our paper makes the following contributions:
\begin{itemize}
    \item \emph{Multimodal instruction-following data}. One key challenge is the lack of vision-language instruction-following data. We present a data reformation perspective and pipeline to convert image-text pairs into an appropriate instruction-following format, using ChatGPT/GPT-4.
    \item \emph{Large multimodal models}. We develop a large multimodal model (LMM), by connecting the open-set visual encoder of CLIP with the language decoder Vicuna, and fine-tuning end-to-end on our generated instructional vision-language data. Our empirical study validates the effectiveness of using generated-purpose instruction-following visual agent. When ensembled with GPT-4, our approach achieves SOTA on the Science QA multimodal reasoning dataset.
    \item \emph{Multimodal instruction-following benchmark}. We present LLaVA-Bench with two challenging benchmarks, with a diverse selection of paired images, instructions and detailed annotations.
    \item \emph{Opne-source}. We release the following assets to the public: the generated multimodal instruction data, the codebase, the model checkpoints, and a visual chat demo.
\end{itemize}


%%%%%%%%%%%%%%%%%%%%%%%%%%%%%%%%%%%%%%%%%%%%%%%%%%%%%%%%%%%%
\section{Related Work}
\noindent\textbf{Multimodal Instruction-following Agents.} In computer vision, existing works that build instruction-following agents can be broadly categorized into two classes: ($\mathrm{i}$) End-to-end trained models, which are separately explored for each specific research topic. For example, the vision-language navigation and take a sequence of actions to complete goals in visual environments. In the image editing domain, given an input image and a written instruction that tells the agents what to do, InstructPix2Pix edits images by following the human instructions. ($\mathrm{ii}$) A system that coordinates various models via LangChain/LLM, such as Visual ChatGPT, X-GPT, MM-REACT, VisProg, and ViperGPT. While sharing the same goal in building instruction-following agents, we focus on developing an end-to-end trained language-vision multimodal model for \emph{multiple} tasks.

\noindent\textbf{Instruction Tuning.} In the natural language processing (NLP) community, to enable LLMs such as GPT-3, T5, PaLM and OPT to follow natural language instructions and complete real-world tasks, researchers have explored methods for LLM instruction-tuning, leading to instruction-tuned counterparts such as InstructGPT/ChatGPT, FLAN-T5, FLAN-PaLM, and OPT-IML, respectively. It turns out that this simple approach can effectively improve the zero and few-shot generalization abilities of LLMs. It is thus matural to borrow the idea from NLP to computer vision. More broadly, the teacher-student distillation ideas with foundation models have been studied  in other topics such as image classification. Flamingo can be viewed as the GPT-3 moment in the multimodal domain, due to its strong performance on zero-shot task transfer and in-context-learning. Other LMMs trained on image-text pairs include BLIP-2, FROMAGe, and KOSMOS-1. PaLM-E is an LMM for enbodied AI. Based on the recent "best" open-source LLM LLaMA, OpenFlamingo and LLaMA-Adapter are open-source multimodal LLMs. While these models present promising task transfer generalization performance in multimodal tasks usually falls short compared to language-only tasks. In this paper, we aim to fill this gap and study its effectiveness. Finally, note that visual instruction tuning is different from visual prompt tuning: the former aims to improve the model's instruction-following abilities, while the latter aims to improve the parameter-efficiency in model adaptation.

%%%%%%%%%%%%%%%%%%%%%%%%%%%%%%%%%%%%%%%%%%%%%%%%%%%%%%%%%%%%
\section{GPT-assisted Visual Instruction data Generation}

%%%%%%%%%%%%%%%%%%%%%%%%%%%%%%%%%%%%%%%%%%%%%%%%%%%%%%%%%%%%
\section{Visual Instruction Tuning}




%%%%%%%%%%%%%%%%%%%%%%%%%%%%%%%%%%%%%%%%%%%%%%%%%%%%%%%%%%%%

\appendix

%%%%%%%%%%%%%%%%%%%%%%%%%%%%%%%%%%%%%%%%%%%%%%%%%%%%%%%%%%%%

\newpage
% \input{checklist.tex}


\end{document}